% Source: https://artofproblemsolving.com/community/c3h1408682_radius_of_circumscribed_circle
% Stable identity: aops-topic-1408682-q01
Let the triangle be $ABC$ with $BC=x$, $CA=y$, $AB=z$, area $K$ and circumcentre $O$.

\textbf{Step 1: $x=2R\sin A$.} Let $BD$ be the diameter through $B$. Then $\angle BCD=90^\circ$ (an angle in a semicircle), and $\angle BDC$ equals $\angle A$ if $D$ and $A$ are on the same arc $BC$, or $180^\circ-\angle A$ otherwise (inscribed angles on the same chord); in both cases $\sin\angle BDC=\sin A$. (If $D=C$, then $\angle A=90^\circ$ and $x=2R$ directly.)

\begin{asy}
usepackage("amsmath");
size(5.5cm);
defaultpen(fontsize(10pt));
path ra(pair P,pair u,pair v,real s=0.2){return P+s*unit(u)--P+s*unit(u)+s*unit(v)--P+s*unit(v);}
pair B=(0,0),C=(14,0),A=(5,12),O=(7,33/8); real R=65/8;
pair Dp=2*O-B;
draw(circle(O,R),gray);
draw(A--B--C--cycle,linewidth(0.9));
draw(B--Dp,blue+dashed); draw(C--Dp,blue+dashed);
draw(ra(C,B-C,Dp-C,0.7),blue);
dot(O); label("$O$",O,SE);
dot(Dp,blue); label("$D$",Dp,NE,blue);
label("$A$",A,N); label("$B$",B,SW); label("$C$",C,SE);
label("$x$",(B+C)/2,S); label("$y$",A+0.3*(C-A),NE); label("$z$",(A+B)/2,W);
label("\mbox{}",(min(currentpicture,user=true)+max(currentpicture,user=true))/2,black); // keeps the text after the figure black when LaTeX typesets the labels inline
\end{asy}

In the right triangle $BCD$, $BC=BD\sin\angle BDC$, so $x=2R\sin A$.

\textbf{Step 2: $R=\dfrac{xyz}{4K}$.} The area is $K=\dfrac{1}{2}yz\sin A=\dfrac{1}{2}yz\cdot\dfrac{x}{2R}=\dfrac{xyz}{4R}$.

\textbf{Step 3: Heron's formula.} $K=\dfrac{1}{4}\sqrt{(x+y+z)(-x+y+z)(x-y+z)(x+y-z)}$. Therefore
\[
\boxed{R=\dfrac{xyz}{\sqrt{(x+y+z)(-x+y+z)(x-y+z)(x+y-z)}}}.
\]

\textbf{Example.} For $13,14,15$: $(42)(16)(14)(12)=112896=336^2$, so $R=\dfrac{13\cdot14\cdot15}{336}=\dfrac{2730}{336}=\boxed{\dfrac{65}{8}}$. (Equivalently $K=84$ and $R=\dfrac{2730}{4\cdot84}$.) In the figure, which shows this triangle with $x=14$, the circumcentre is $O=\left(7,\dfrac{33}{8}\right)$ when $B=(0,0)$ and $C=(14,0)$.
